\documentclass{gescons}

\genre {Entrevista}
\author{Daniel Mamede}
\title{Ações Antifanatismo: Intensidade Lúcida em Autopesquisas e Autorreciclagens}
\paginaurl{https://www.youtube.com/live/8pyO5IIiHRI}

\begin{document}
    \makeentrevistatitle

    \coverart{back/Daniel-Mamede}
    \begin{multicols}{2}

\begin{center}
    \includegraphics[width=6cm]{articles/entrevista/mockups/Daniel-Mamede.png}
\end{center}

%\noindent\includegraphics[width=9cm, height=10cm]{example-image} 

\subsubsection{1. Qual foi a motivação para a escrita da obra? Por que a definição deste tema para publicação de um livro?}

O livro “Ações Antifanatismo” representa muito para mim, por ter sido o resumo dos meus 10 primeiros anos de contato e vivência do paradigma consciencial. 

Escrito, em princípio, para mim mesmo, a fim de organizar, reeducar e curar aspectos negativos das automanifestações, muitas delas tendentes ou vinculadas ao fanatismo, o livro mostrou \textit{a vida teática no papel,} isto é, materializou, grafou e vincou a autopesquisa conscienciológica sendo vivenciada dia a dia, pensene a pensene, aprendizado a aprendizado, resultando em material com enorme potencial para ajudar outras pessoas que possam estar em apuros com o mesmo assunto, notadamente os intermissivistas, os quais partilham de propósitos e objetivos similares aos meus para esta vida humana.



\subsubsection{2. Quais foram as principais percepções, intra e extrafísicas, durante a pesquisa e a escrita da obra? E posterior ao lançamento?}

O que este livro mais tem são percepções e parapercepções vivenciadas durante os processos de reciclagem, pois seu estilo de escrita predominante foi o do “\textit{in media res}”, isto é, as experiências da vida eram grafadas e autocriticadas de modo praticamente concomitante às ocorrências. Vale muito a pena conferir. 


\subsubsection{3.       Qual o maior aprendizado com a escrita desta obra?}

Nos últimos 10 anos, nunca deixei passar mais que 24 horas entre o intervalo da ocorrência dos fatos ou parafatos e as devidas anotações. Por isso, o efeito das vivências descritas no livro traz as energias dos momentos das aprendizagens disruptivas de modo bastante intenso e real. E essa rotina continua firme na minha vida, mesmo após a conclusão do livro, pois a autopesquisa conscienciológica nunca pode parar. Este foi meu maior aprendizado ao ver o livro concluído.

\begin{pullquote}
``O efeito das vivências descritas no livro traz as energias dos momentos das aprendizagens disruptivas de modo bastante intenso e real.''
\end{pullquote}


\subsubsection{4.       O que poderia dizer como incentivo para que mais pesquisadores invistam na publicação de obras conscienciológicas?}

Fazer autopesquisas diárias e comprometidas como as descritas em “Ações Antifanatismo” é um dos caminhos mais tranquilos para quem quiser publicar algum livro. Ele vai acontecer naturalmente. Vai estar pronto antes mesmo de o autor o conceber. Foi assim comigo. O acúmulo dos pensenes materializados no papel gerou banco de dados rico e vasto para eu poder explorá-lo da maneira que entendesse por bem, e assim o fiz justamente para trabalhar os traços de temperamento que entendi serem os mais sérios e prioritários em meus processos de reciclagem, no caso, vinculados às posturas fanáticas. Aí, foi só organizar. A rigor, o que direcionou e conduziu a escrita do livro não foi a atenção ao traf\textit{a}r do fanatismo, mas sim ao traf\textit{o}r da valorização e do comprometimento com a ininterruptibilidade da autopesquisa conscienciológica.


    
    
    \end{multicols}
\end{document}
